\label{chap:p3-83-27-menos_un_doceavo}

\section{Dominio Hensel--Witt, mapas de Teichmüller y objetivo racional \(-1/12\)}

El NTHW es el paquete de datos y mapas
\[
\NTHW=
\bigl(
\mathscr R_\pi^{\rm micro},
\Gnine,
\Cdoce,
K,
\mathcal W_{12\to24},
\mathcal I_{6\to8}
\bigr),
\]
con las compatibilidades siguientes:
\begin{enumerate}
  \item las regiones Hensel comparten una palabra cociente y conservan
  distinta memoria;
  \item \(\Gnine\) transporta la orientación y produce monodromía;
  \item \(\Cdoce\) organiza el observable firmado y el sello \(K\);
  \item la lectura Witt convierte soportes del mismo borde en hexadas y
  estrellas;
  \item la extensión \(W_{12}\to W_{24}\) produce la interfaz
  hexada--octada;
  \item los lectores arquimediano, excepcional y dimensional reciben
  invariantes del mismo estado.
\end{enumerate}

\begin{figure}[!htbp]
\centering
\resizebox{\textwidth}{!}{\begin{tikzpicture}[font=\small,node distance=1.2cm and 1cm,
  nodebox/.style={draw,rounded corners=2mm,align=center,minimum width=3cm,minimum height=1.05cm},
  arr/.style={-{Stealth[length=2.4mm]},thick}]
  \node[nodebox,draw=hmtblue,fill=hmtlightblue] (hensel) at (-5,2.3) {cilindros henselianos\\cociclo de acarreo y \(\MonNine\)};
  \node[nodebox,draw=hmtblue,fill=hmtlightblue] (regions) at (-5,-.1) {cinco regiones de \(\pi\)\\regiones de \(e,\varphi\)};
  \node[nodebox,draw=hmtgold,fill=hmtlightgold,minimum width=3.8cm,minimum height=2.05cm,font=\bfseries] (core) at (0,1.1) {\(\mathscr T_{\rm HW}\)\\Correspondencias\\Hensel--Witt};
  \node[nodebox,draw=hmtviolet,fill=hmtlightviolet] (c12) at (0,-1.8) {\(\Cdoce\), \(\Usgn\), \(K\)\\diferencias \(D_3,D_4,Q\)};
  \node[nodebox,draw=hmtgreen,fill=hmtlightgreen] (witt) at (5,2.3) {Witt \(W_{12}\)\\soportes y \(M_{12}\)};
  \node[nodebox,draw=hmtgreen,fill=hmtlightgreen] (oct) at (5,-.1) {inclusión de incidencia \(6\to8\)\\canales \(90/120\), normalización \(-1/12\)};
  \node[nodebox,draw=hmtred,fill=hmtlightred] (readers) at (5,-2.5) {lectores coordinados\\valor y simetría};
  \draw[arr,hmtblue] (hensel)--(core);
  \draw[arr,hmtblue] (regions)--(core);
  \draw[arr,hmtgold] (core)--(c12);
  \draw[arr,hmtgreen] (core)--(witt);
  \draw[arr,hmtgreen] (core)--(oct);
  \draw[arr,hmtred] (c12)--(readers);
  \draw[arr,hmtred] (witt.east)--++(1.9,0)|-(readers.east);
  \draw[arr,hmtred] (oct)--(readers);
  \node[align=center,hmtgray,font=\footnotesize] at (0,-3.45) {Compatibilidad de cilindros, dodecafase e incidencias sobre un mismo estado enriquecido.};
\end{tikzpicture}
}
\caption{El NTHW como interfaz tipada. La tarea matemática central es hacer
conmutar el diagrama, no acumular coincidencias alrededor de un nombre.}
\label{p3-83-c27-s1:fig:nthw}
\end{figure}

El NTHW da un nombre formal al lugar que antes se describía como «Piedra
Rosetta». La sustitución es conceptual: una piedra traduce por semejanza; una
transducción especifica entrada, salida, dato conservado y pérdida de
información.

\section{Aplicación de Teichmüller del código ternario al anillo de Witt}

La frontera reversible de APP selecciona, con calibres declarados, una matriz
\(A_W\) sobre \(\FF_3\) que satisface
\[
A_WA_W^T=2I\pmod3.
\]
El código gráfico
\[
C_W=\{(q,qA_W):q\in\FF_3^6\}
\]
es el Golay ternario extendido. Su enumerador de pesos es
\[
1+264y^6+440y^9+24y^{12}.
\]
Las \(264\) palabras de peso seis, identificadas por signo, dan \(132\)
soportes. Cada subconjunto de cinco puntos pertenece a una única hexada:
\[
W_{12}=S(5,6,12).
\]

El campo de \(495\) involuciones de estrella genera un grupo de orden
\(95\,040\), reconocido como \(M_{12}\). Una estrella aislada no genera el
grupo. El marco HMT ordenado tiene estabilizador trivial y por ello fija un
calibre dentro de la acción.

\section{Incidencia hexada--octada y canales modulares \(90/120\)}

Sea \(H\) una hexada y \(O_H\) una octada que la contiene. Al marcar un punto
y un par de puntos de la hexada se obtienen
\[
F_6(H)=H\times\binom H2,
\qquad
|F_6|=6\binom62=90.
\]
Si el punto marcado puede recorrer la octada:
\[
F_8(H)=O_H\times\binom H2,
\qquad
|F_8|=8\binom62=120.
\]
La inclusión \(H\hookrightarrow O_H\) induce \(F_6\hookrightarrow F_8\).
Su defecto normalizado es
\[
\boxed{
\frac{90-120}{24\binom62}
=\frac{6-8}{24}
=-\frac1{12}.}
\]
Este teorema combinatorio es anterior a la lectura angular de Barbero. Los
dos objetos comparten la interfaz \(6\to8\), pero no son idénticos.

\section{4 rutas compatibles hacia \texorpdfstring{\(-1/12\)}{-1/12}}

La construcción contiene cuatro realizaciones que deben exponerse juntas y no
contarse como cuatro pruebas independientes de una misma procedencia.

\subsection{Realización angular mediante el incremento elemental dodecafásico}

Un diente de \(\Cdoce\) mide \(30^\circ\). Entonces
\[
\frac{30}{120}-\frac{30}{90}
=\frac14-\frac13
=-\frac1{12}.
\]
Es la sombra angular del patrón \(90/120\).

\subsection{Pseudoinversa orientada del operador de Gram sobre el subespacio superviviente}

En el salto certificado \(60\to81\), dos fibras tienen tamaños \(0\) y
\(12\). El operador de Gram es
\[
K_{60,81}=\operatorname{diag}(0,12).
\]
Sobre el subespacio vivo, su pseudoinverso vale \(1/12\); con la orientación
negativa de borde,
\[
E_{\rm surv}=-K_{60,81}^{+}
\]
produce \(-1/12\).

\subsection{Extracción de escala por 2.ª diferencia triádica}

La realización mediante un extractor de escala triádico es la construcción por
segunda diferencia desarrollada íntegramente en
\cref{sec:c27-segunda-diferencia-triadica}. Esta primera comparecencia fija su
función dentro del inventario de realizaciones; la residencia indicada conserva
las definiciones de \(T_L\), \(a(L)\), \(A_1(L)\), \(C(L)\), la independencia
del residuo respecto de \(L\) y el estatuto exacto de la normalización.

\subsection{Realización modular mediante el cociente \(\eta(\tau)/\eta(3\tau)\) y la función \(t_3\)}

La función eta satisface
\[
\frac{\eta(\tau)}{\eta(3\tau)}
=q^{-1/12}\prod_{n\ge1}(1+q^n+q^{2n})^{-1}.
\]
El exponente procede de \(1/24-3/24=-1/12\). Al tomar doce copias:
\[
t_3(q)=\left(\frac{\eta(\tau)}{\eta(3\tau)}\right)^{12}
=q^{-1}-12+54q-76q^2-243q^3+\cdots.
\]
El coeficiente \(54\) coincide con el defecto APP. El corredor
\[
j=\frac{(t_3+27)(t_3+243)^3}{t_3^3}
\]
conecta después con la rama modular clásica.

\begin{remark}[Qué se ha demostrado en el corredor]
Son exactos el defecto APP \(54\), el exponente eta \(-1/12\), el defecto
hexada--octada y las identidades modulares una vez fijado \(t_3\). La
afirmación fuerte es que todos sean imágenes naturales de un único operador
TPK. Esa naturalidad debe demostrarse mediante el diagrama del NTHW; no se
deduce sólo de repetir los mismos enteros.
\end{remark}


\section{5 realizaciones algebraicas de \texorpdfstring{\(-1/12\)}{-1/12} y separación de sus dominios}

Además del defecto normalizado de la inclusión hexada--octada, aparecen
cinco realizaciones operatoriales del mismo racional. Sus dominios se
especifican por separado; la igualdad de valores no identifica los
operadores sin un entrelazador adicional.

\subsection{Diferencia de periodicidades dodecafásicas}

Un diente de \(\Cdoce\) mide \(30^\circ\). Entonces
\[
\frac{30}{120}-\frac{30}{90}
=\frac14-\frac13
=-\frac1{12}.
\]
Es la diferencia orientada entre los periodos de \(90^\circ\) y
\(120^\circ\) en esta realización.

\subsection{Pseudoinversa del operador de Gram sobre la componente viva}

En el salto exacto \(60\to81\), dos fibras tienen tamaños \(0\) y
\(12\). El operador de Gram es
\[
K_{60,81}=\operatorname{diag}(0,12).
\]
Sobre la componente no nula de su rango, la pseudoinversa vale \(1/12\); con la orientación
negativa de borde,
\[
E_{\rm surv}=-K_{60,81}^{+}
\]
produce \(-1/12\).

\subsection{2.ª diferencia entre escalas triádicas}
\label{sec:c27-segunda-diferencia-triadica}

Definamos
\[
T_L=\sum_{n=1}^{L-1}n(L-n)=\frac{L(L^2-1)}6,
\qquad
a(L)=-\frac{T_L}{L^2}=-\frac L6+\frac1{6L}.
\]
La primera resta de escala es
\[
A_1(L)=a(L)-\frac{a(3L)}3=\frac4{27L},
\]
y la segunda
\[
C(L)=LA_1(L)-\frac{3L}{3}A_1(3L)=\frac8{81}.
\]
El residuo \(8/81\) es independiente de \(L\). Para obtener
\(-1/12\) se aplica la normalización declarada
\[
R(L)=-\frac{27}{32}C(L)=-\frac1{12}.
\]
El cálculo deriva \(8/81\); no deriva por sí solo el factor \(-27/32\).

\subsection{Cadena modular}

La función eta satisface
\[
\frac{\eta(\tau)}{\eta(3\tau)}
=q^{-1/12}\prod_{n\ge1}(1+q^n+q^{2n})^{-1}.
\]
El exponente procede de \(1/24-3/24=-1/12\). Al tomar doce copias:
\[
t_3(q)=\left(\frac{\eta(\tau)}{\eta(3\tau)}\right)^{12}
=q^{-1}-12+54q-76q^2-243q^3+\cdots.
\]
El coeficiente \(54\) coincide con el defecto APP\@. La identidad racional
\[
j=\frac{(t_3+27)(t_3+243)^3}{t_3^3}
\]
establece la conexión con la rama modular clásica.

\Needspace{32\baselineskip}
\subsection{Defecto de los proyectores de diferencia cíclica}

Sea \(S\) el desplazamiento cíclico sobre \(\QQ^{12}\) y definamos
\[
 D_r=I-S^r.
\]
El núcleo de \(D_r\) está formado por los vectores constantes en cada órbita
de \(S^r\); por tanto,
\[
 \dim\ker D_r=\gcd(12,r).
\]
Sean \(\Pi_{\ker D_3}\) y \(\Pi_{\ker D_4}\) los proyectores racionales sobre
\(\ker D_3\) y \(\ker D_4\), y sea
\(\tau(T)=\operatorname{Tr}(T)/12\) la traza normalizada. Entonces
\[
 \boxed{
 \tau(\Pi_{\ker D_3}-\Pi_{\ker D_4})
 =\frac{\operatorname{rank}\Pi_{\ker D_3}
       -\operatorname{rank}\Pi_{\ker D_4}}{12}
 =\frac{3-4}{12}
 =-\frac1{12}.}
\]
Esta realización expresa el racional como defecto transversal de los pasos
tres y cuatro en el ciclo de doce fases.

\begin{remark}[Dominios de las realizaciones]
Son exactos el defecto APP \(54\), el exponente eta \(-1/12\), el defecto
combinatorio de la elevación hexada--octada, y las
identidades modulares una vez fijado \(t_3\). También son exactos el valor
de la pseudoinversa, el extractor normalizado y el defecto de proyectores.
Cada igualdad conserva su dominio y su aplicación correspondiente.
\end{remark}
